% !TeX root = ../main.tex
\section{Auditing the Judge: Gain Retention}
\label{sec:validity}

Everything so far was graded by the model that supplied the training reward and that also
played the patient. To test whether the conclusions are properties of the therapist or of
that grader, we re-scored the \emph{entire} grid --- every conversation, every rubric, every
model state, $22{,}272$ cells --- with \heldoutjudge{}: a different model family that never
played the patient and never touched training.

\paragraph{The measuring instruments are repeatable.}
Reproducibility is measured rather than assumed. Re-scoring the same conversations with
only the sampling seed changed gives ICC(2,1) \citep{shrout1979icc} of $0.98$--$0.99$ (Q1),
$0.96$--$0.99$ (Q2) and $0.86$--$0.94$ (\mici{}) for the primary oracle, and
$0.95$--$0.98$ / $0.94$--$0.96$ / $0.53$--$0.93$ for the held-out judge
(Appendix~\ref{app:reliability}, Table~\ref{tab:icc}) --- ``excellent'' on the reward
instruments by the \citet{koo2016icc} guideline. This matters for reading the next
paragraph: two noisy raters cannot agree perfectly even if they measure the same thing, so
cross-judge agreement must be read against the attenuation ceiling
$\sqrt{\mathrm{ICC}_a \cdot \mathrm{ICC}_b}$, never against $1.0$
(Table~\ref{tab:ceiling}).

\paragraph{The rankings survive.}
Across all $1{,}848$ pairwise (arm $\times$ metric) contrasts in this comparison, $88.3\%$
keep their sign under the held-out judge, rising monotonically with effect size to $94.1\%$
at $|\Delta|\!\ge\!0.10$, $97.0\%$ at $\ge\!0.25$ and $98.9\%$ at $\ge\!0.50$ (Appendix~\ref{app:signs},
Table~\ref{tab:signladder}): the two
graders disagree only about differences too small to claim. All 18 contrasts this paper
leans on --- including the steelman PTO@10 $-$ GRPO@8 and the regression claim
GRPO@8 $-$ GRPO@10 --- are preserved with bootstrap CIs excluding zero, and the held-out
judge in fact \emph{widens} the headline Q1 gap ($+0.77$ vs the primary's $+0.53$). A
two-way decomposition of the arm means puts only $1.2$--$6.9\%$ of variance in the
arm~$\times$~judge interaction, the only component that could invalidate an ordering; the
judges differ on \emph{level} (the held-out judge is $1.2$--$1.7$ points harsher), which
cancels in every contrast we report.

\paragraph{One rubric genuinely disagrees.}
MITI is the exception and we flag it as a limitation rather than a footnote: only $3.6\%$
of its arm-mean variance is between-arm signal, a single-judge MITI ranking is $0.65$
dependable against ${\approx}0.90$ for every other rubric, and it preserves its sign on only
$77.5\%$ of contrasts --- and, unlike every other rubric, still only $88.2\%$ at
$|\Delta|\!\ge\!0.25$, i.e.\ it can flip at a gap large enough to want to report.
\textbf{MITI arm differences in this paper should be treated as provisional unless both
graders agree on the direction.} \mici{} agrees only weakly per conversation
($r=0.20$--$0.55$); against the measured ceiling it recovers ${\approx}39\%$ of achievable
agreement where Q1/Q2 recover ${\approx}85\%$, so the residual is mostly genuine construct
disagreement about what counts as MI-inconsistent. We therefore state the sycophancy claim
at the \emph{contrast} level (both graders agree GRPO@10 is more MI-inconsistent than
PTO@10, and that \mici{} rises from base in both arms) and not as a precise rate.

\paragraph{Gain retention separates the arms.}
The sharpest use of a held-out grader is not agreement but \emph{retention}. Because the
primary oracle was the training reward and the second judge never touched training,
\[
  \text{retention} \;=\;
  \frac{\Delta_{\text{held-out judge}}}{\Delta_{\text{trained-against judge}}}
\]
is a train/test generalization ratio for each arm's gain over base: ${\approx}1$ means a
behaviour change both graders see, ${\approx}0$ means a gain that existed only inside the
optimized grader. On Q1 the arms separate and their intervals do not overlap:
PTO@10 $0.80\,[0.68, 0.93]$ versus GRPO@10 $0.28\,[0.06, 0.43]$ (GRPO@8:
$0.64\,[0.53, 0.78]$). Under a judge that never graded during training, GRPO's net
ten-iteration Q1 gain is ${\approx}0.19$ points --- close to where it started --- against
the ${\approx}0.68$ the primary grader credits it. Q2 retention, by contrast, is flat at
$0.80$--$0.85$ with fully overlapping intervals, which is ordinary scale compression
between graders and is exactly the null result that makes the Q1 divergence readable as
arm-specific rather than as a grader artifact.

\begin{table}[t]
  \centering\small
  \setlength{\tabcolsep}{3.4pt}
  \begin{tabular}{@{}lcccccccccc@{}}
    \toprule
    iter & 1 & 2 & 3 & 4 & 5 & 6 & 7 & 8 & 9 & 10\\
    \midrule
    PTO  & .97 & .84 & .89 & .94 & .98 & .97 & .94 & .89 & .88 & \textbf{.80}\\
    GRPO & 1.13 & .79 & .89 & .79 & .73 & .57 & .70 & .64 & .03 & \textbf{.28}\\
    \bottomrule
  \end{tabular}
  \caption{Q1 gain retention by iteration. The arms are indistinguishable for three
  iterations, then separate: PTO holds $0.80$--$0.98$ for the whole run while GRPO decays
  monotonically in trend. This is an onset curve, not an endpoint accident.}
  \label{tab:retention}
\end{table}

\begin{figure}[t]
  \centering
  \includegraphics[width=\linewidth]{figures/multijudge_retention_trajectory.png}
  \caption{Gain retention across iterations (Table~\ref{tab:retention}, plotted). The
  held-out grader stops crediting GRPO's gains \emph{progressively} --- the signature of a
  policy drifting onto grader-specific features rather than of a bad final checkpoint.}
  \label{fig:retention}
\end{figure}

\paragraph{Why this is the load-bearing evidence.}
Retention is measured on Q1 --- an instrument \emph{inside} the training reward --- which
makes it immune to the objection that outside-the-reward instruments (\mici{}, MITI) might
simply be measuring something the training never targeted. It requires no human coding, no
new conversations, and one extra grading pass. And it is an onset curve
(Table~\ref{tab:retention}, Figure~\ref{fig:retention}): the divergence begins at
iteration~3 and grows, which is what optimizer drift looks like and is not what a single
bad checkpoint looks like. We propose it as a default diagnostic for any system trained
against an LLM judge: \emph{report how much of your gain a grader that never saw training
still credits, per iteration.}
